% Reversed code order and high-mask quotient test: 111 arithmetic instructions.
\section{Reversing the packed codes}\label{sec:reversed-packing}

The inequality $e+\ell\mathcal C^2<q$ costs a multiplication and two
additions, including its positive slack. The already computed packed
quantity can replace the product, but its two components must still be
recovered uniquely. Reversing their order and shortening the coefficient
code makes the third mask enforce that uniqueness. The resulting bound
costs two additions.

\subsection{The revised fixed index and system}

Retain the quadratic layout, the polynomial $\mathcal D$, and the mask
$\ell_0$ from Section~\ref{sec:quadratic-masks}. Put $K=t_*+1$ and keep
$L=5^{16}$. The fixed numerical layout satisfies
\begin{equation}\label{eq:reverse-margin}
 L>3K+2.
\end{equation}
Shorten the dense part of the coefficient polynomial and reverse the
order of the two codes:
\begin{equation}\label{eq:reverse-code}
 e_0(B)=z\sum_{j=0}^{K-1}B^j+\mathcal D(B),\qquad
 V=\ell_0(Z)+e_0(Z)Z^L.
\end{equation}
Both $e_0$ and $\ell_0$ now have degree at most $t_*$; their coefficients
are still in $[0,Z)$, with the coefficients of $e_0$ positive.
Choose the fixed power of two $H$ so that
\begin{equation}\label{eq:reverse-index}
 H>\max\{2Z^{2L+1},\ 2^{t_*+4}Z(1890)^2,\ 3L,\ 16\}.
\end{equation}
This includes $H>4Z(1890)^2$ and all earlier radix bounds. These index
choices depend only on the represented r.e.\ set.

Start from the system of Theorem~\ref{thm:best112}, retaining
$B=Hb^2$, $\mathcal C=1+xB+g$, and all its Pell equations. The changed
coding equations and packing expressions are
\begin{align}
 Y&=\ell+eq,&
 Y+\mathcal C^2+\alpha&=q^2,&
 Y&=V+t\theta,\label{eq:reverse-bound}\\
 S_3&=(2e-Z\lambda)\mathcal C^2+B\lambda(1+q),&
 \sigma&=S_3,\label{eq:reverse-sigma}\\
 S&=g+q^2(Y+q^2S_3),\label{eq:reverse-S}\\
 T+1&=q^2(1+\theta\lambda)-(b-1)\ell
                    +(B-2)\ell q^4,\label{eq:reverse-T}\\
 r&=S(n^2-n)+(T+1)(n^2-1).\label{eq:reverse-r}
\end{align}
Here $Y,S_3,S,T$ remain computed abbreviations. The only new unknown is
$\sigma$, required to be positive. Its equality to the already computed
$S_3$ is a free equality test. The unchanged coding equations include
\[
 b=x+\beta,\qquad \lambda+q^2=1+\lambda B,\qquad
 \theta+Z=B,\qquad n=q^8.
\]

\subsection{Bounds before decoding and a possible borrow}

Equation~\eqref{eq:reverse-bound} implies
\[
 \ell<Y<q^2,\qquad e<q,\qquad \mathcal C<q,\qquad g<q.
\]
In particular $q>B$, since $x>0$ and $\mathcal C>B$. With
$N=\theta\lambda$, the geometric equation and the radix bounds give
\[
 q<\lambda<2q^2/B,\qquad b<N<q^2-1.
\]
The new unknown supplies $S_3>0$. Before any decoding we also have
\[
 S_3<2q^3+B\lambda(1+q)<5q^3<q^4.
\]
The three blocks of $S$ therefore have widths $q^2,q^2,q^4$ and give
$0<S<n=q^8$.

The first raw mask $q^2-1-(b-1)\ell$ may be negative at this stage.
Its complete packed value is positive, because
\[
 q^2(1+N)-(b-1)\ell>q^2,
 \qquad (B-2)\ell<q^3.
\]
In addition $T+1<q^7+q^4+q^2<q^8$. Thus $0\le T<n$ and
$r\ge n^2-1\ge n$. The hypotheses used by the unchanged Pell
implications hold without assuming the sign of that first raw mask.
Those implications give $q=B^L$, the required powers of two, and
$\tau_2(S,T)=0$.

To read the masks, write
\[
 d=\left\lfloor\frac{(b-1)\ell}{q^2}\right\rfloor,
 \qquad u=(b-1)\ell-dq^2.
\]
Then $0\le d<b$, $0\le u<q^2$, and the actual block decomposition is
\[
 T=(q^2-1-u)+q^2(N-d)+q^4(B-2)\ell.
\]
All three blocks are in their stated ranges. Every base-$B$ digit of
$N=(B-Z)\lambda$ equals $B-Z$, in positions $0,\ldots,2L-1$.
Since $d<b<B-Z$, subtracting $d$ changes only the unit digit. The middle
no-carry condition consequently makes every nonunit digit of $Y$ less
than $Z$, and makes its unit digit at most $Z+d-1$. Hence
\[
 Y<\frac{Zq^2}{B-1}+b,
 \qquad
 (b-1)\ell<\frac{(b-1)Zq^2}{B-1}+b(b-1)<q^2.
\]
For the last inequality, \eqref{eq:reverse-index} gives
$B-1>4(b-1)Z$ and $q>4b$. Therefore $d=0$ after all.

All digits of $Y$ are now less than $Z$. The digit-evaluation lemma
applied to \eqref{eq:reverse-code} and the packed congruence gives
\begin{equation}\label{eq:reverse-quotient}
 Y=\ell_0(B)+e_0(B)q,\qquad
 \ell=\ell_0(B)+mq,\qquad e=e_0(B)-m
\end{equation}
for an integer $m$ with $0\le m<e_0(B)$. Positivity of $\ell$ and
$\ell_0(B)<q$ exclude negative $m$. The following argument rules out
positive $m$ as well.

\subsection{The high mask checks the quotient}

The fixed support gives $(b-1)\ell_0(B)<q$. The low $q$ block of the
first mask is therefore $q-1-(b-1)\ell_0(B)$. Since $g<q$, its no-carry
condition decodes exactly the earlier low coordinate positions, with
every digit less than $b$. Extra coordinates at positions belonging to
$mq$ cannot occur in $g$, because those positions are at least $L$.
Thus $\mathcal C$ has degree at most $t_*$, constant digit $1$, and input
digit $x$. Writing $A$ for the sum of the coefficients of
$\mathcal C^2$, we obtain
\begin{equation}\label{eq:reverse-A}
 A=\mathcal C(1)^2<(1890b)^2,
 \qquad 4\le ZA<B/4.
\end{equation}
Here $\mathcal C(1)$ denotes evaluation of its digit polynomial, not of
the integer expression $1+xB+g$ with $g$ held constant.

By \eqref{eq:reverse-quotient}, both $e$ and $m$ are smaller than
$e_0(B)<B^K$. The decoded $\mathcal C$ also satisfies $\mathcal C<B^K$.
The explicit margin \eqref{eq:reverse-margin} consequently gives an
integer bound which includes all carries:
\begin{equation}\label{eq:reverse-small-product}
 2e\mathcal C^2<2B^{3K}<B^L=q.
\end{equation}

Put
\[
 X=(B-2)m=\sum_{j=0}^{t_*+1}X_jB^j,
 \qquad 0\le X_j<B.
\]
Since $(B-2)\ell_0(B)<q$, the digits of the third mask at positions
$L,\ldots,L+t_*+1$ are exactly the $X_j$. In this window the coefficient
of $\lambda\mathcal C^2$ is always $A$: the window is above $2t_*$
and below $2L$. By \eqref{eq:reverse-small-product}, the positive integer
$2e\mathcal C^2$ has no digit there. The offset
$B\lambda(1+q)$ has raw coefficient $1$ at position $L$ and raw
coefficient $2$ at the subsequent positions of the window.

The incoming carry at $L$ lies in $\{-1,0,1\}$. Indeed the raw negative
coefficients below $L$ have absolute value at most $ZA<B/4$, while the
positive contribution below $L$ is less than $q+q/(B-1)<2q$.
Thus the normalized digit at $L$ lies between $B-ZA$ and $B-ZA+2$;
because $ZA\ge4$, its outgoing carry is $-1$. The later digits in the
window are $B-ZA+1$ and their outgoing carry remains $-1$. In particular
every digit in this window is at least $B-ZA$.

No carry with the third mask now implies $X_j\le ZA-1$. Consider the
ordinary polynomial $F(U)=\sum_jX_jU^j$. Its value
$F(B)=(B-2)m$ is divisible by $B-2$, so $F(2)$ is divisible by $B-2$.
But \eqref{eq:reverse-index} and \eqref{eq:reverse-A} give
\[
 0\le F(2)<ZA\,2^{t_*+2}
       <Z(1890)^2b^2 2^{t_*+2}<B/4<B-2.
\]
Therefore $F(2)=0$. All coefficients of $F$ are nonnegative, so $X=0$
and $m=0$. We have recovered both canonical codes
$e=e_0(B)$ and $\ell=\ell_0(B)$.

Every target position is below $K$. Below that cutoff,
$e_0-z\lambda=\mathcal D$, and the offset in $S_3/2$ still supplies
the digit $B/2$ at every target. The coefficient bound and the support
separation of Section~\ref{sec:quadratic-masks} apply unchanged.
Carries from overlapping offset blocks above $L$ cannot affect these
lower positions. Thus all original quadratic residuals vanish.

\subsection{Positive witnesses and the verified count}

Conversely, encode a genuine quadratic-system solution, assign one dummy
coordinate the value $1$, and choose a power of two $b$ larger than $x$
and all coordinate values. Use the canonical $e_0(B)$ and $\ell_0(B)$.
Their degrees are at most $t_*$, and \eqref{eq:reverse-margin} gives
$Y+\mathcal C^2<q^2$, with a strictly positive slack.
It also gives $Z\mathcal C^2<q$, so
\[
 S_3\ge\lambda\bigl(B(1+q)-Z\mathcal C^2\bigr)>0.
\]
Set $\sigma=S_3$. The quotient in the packed congruence is positive
because $B>Z$ and its fixed digit polynomial has positive nonconstant
coefficients. The three no-carry conditions hold, and the preceding
Pell necessity constructions supply all remaining positive unknowns.

\begin{theorem}\label{thm:best111}
For each r.e.\ set, an effective admissible index $(Z,V,H)$ defined by
\eqref{eq:reverse-code}--\eqref{eq:reverse-index} gives a system in
33 positive unknowns and 21 equations such that, for positive input $x$,
membership is equivalent to a valid certificate with
\textbf{111 arithmetic instructions}: 61 multiplications and 50 additions.
Equality tests and supplied parameters are free. Generating the fixed
numerals $2,4,5,L$ from $1$ adds seven instructions, giving 118 under
that convention.
\end{theorem}
\begin{proof}
The equivalence, including both positive-witness constructions, was
proved above. Relative to Theorem~\ref{thm:best112}, the product
$\ell\mathcal C^2$, addition of $e$, and addition of the positive slack
are replaced by $Y+\mathcal C^2$ and addition of that slack, saving one
multiplication. The reversed code $\ell+eq$ has the same cost as
$e+\ell q$. Computing $S$ still uses two multiplications and two
additions. The $T+1$ calculation uses
\[
 q^2(\lambda+q^2-Z\lambda)
       -(b-1)\ell+(B-2)(\ell q^4),
\]
which has the same cost as before. Replacing $1+q^2$ by $1+q$ also has
equal cost, and $\sigma=S_3$ is an equality test on an existing computed
register. The exact primitive instructions, equality tests, and
triangular polynomial residual corrections are checked by
\file{round7\_1980\_certificate.py}. The numeral construction is unchanged.
\end{proof}
